\documentclass[10pt, aspectratio=169]{beamer}

% Theme Selection
\usetheme{Madrid}
\usecolortheme{seahorse}

% Packages
\usepackage{amsmath, amssymb, amsfonts}
\usepackage{graphicx}
\usepackage{booktabs}
\usepackage{url}

% Presentation Metadata
\title[Linear Algebra Fundamentals]{Linear Algebra: Vectors, Matrices, and Transformations}
\subtitle{A Geometric and Algebraic Perspective}
\author{Alen Francis Joseph}
\institute{M.Tech CSE (DS and AI), DCS, CUSAT}
\date{\today}

\begin{document}

% -----------------------------------------------------------------------------
\begin{frame}
    \titlepage
\end{frame}

\begin{frame}{Presentation Outline}
    \tableofcontents
\end{frame}

% =============================================================================
\section{Vectors and Coordinate Systems}
% =============================================================================

\begin{frame}{Vectors: Addition and Scalar Multiplication}
    \begin{definition}[Vectors and Coordinates]
        A \textbf{vector} is a geometric object with magnitude and direction, represented in a coordinate system as an array of numbers (components).
    \end{definition}

    \begin{block}{Derivation / Rules}
        Vector addition is performed element-wise. Scalar multiplication distributes the scalar to each individual component:
        $$\begin{bmatrix} x_1 \\ y_1 \end{bmatrix} \pm \begin{bmatrix} x_2 \\ y_2 \end{bmatrix} = \begin{bmatrix} x_1 \pm x_2 \\ y_1 \pm y_2 \end{bmatrix} \quad \text{and} \quad c \begin{bmatrix} x \\ y \end{bmatrix} = \begin{bmatrix} cx \\ cy \end{bmatrix}$$
    \end{block}
\end{frame}

\begin{frame}{Vectors: Addition and Scalar Multiplication (Example)}
    \begin{exampleblock}{Example Problem}
        \textbf{Problem:} Let $\vec{u} = \begin{bmatrix} 2 \\ -1 \end{bmatrix}$ and $\vec{v} = \begin{bmatrix} 3 \\ 4 \end{bmatrix}$. Find $2\vec{u} + \vec{v}$. \\
        \vspace{0.5cm}
        \textbf{Solution:} 
        $$2\vec{u} + \vec{v} = 2\begin{bmatrix} 2 \\ -1 \end{bmatrix} + \begin{bmatrix} 3 \\ 4 \end{bmatrix}$$
        $$= \begin{bmatrix} 4 \\ -2 \end{bmatrix} + \begin{bmatrix} 3 \\ 4 \end{bmatrix}$$
        $$= \begin{bmatrix} 4 + 3 \\ -2 + 4 \end{bmatrix} = \begin{bmatrix} 7 \\ 2 \end{bmatrix}$$
    \end{exampleblock}
\end{frame}

\begin{frame}{Linear Combinations, Span, and Basis Vectors}
    \begin{definition}[Linear Combination and Basis]
        A \textbf{linear combination} is the sum of scaled vectors: $c_1\vec{v}_1 + \dots + c_n\vec{v}_n$. The \textbf{span} is the set of all possible linear combinations. \textbf{Basis vectors} are a linearly independent set that spans the space.
    \end{definition}

    \begin{block}{Derivation / Concept}
        The standard 2D basis vectors are $\hat{i} = \begin{bmatrix} 1 \\ 0 \end{bmatrix}$ and $\hat{j} = \begin{bmatrix} 0 \\ 1 \end{bmatrix}$. Any vector $\vec{w}$ in $\mathbb{R}^2$ can be uniquely derived as $x\hat{i} + y\hat{j}$.
    \end{block}
\end{frame}

\begin{frame}{Linear Combinations, Span, and Basis Vectors (Example)}
    \begin{exampleblock}{Example Problem}
        \textbf{Problem:} Express the vector $\vec{w} = \begin{bmatrix} 5 \\ -7 \end{bmatrix}$ as a linear combination of the standard basis vectors $\hat{i}$ and $\hat{j}$. \\
        \vspace{0.5cm}
        \textbf{Solution:}
        $$\vec{w} = \begin{bmatrix} 5 \\ 0 \end{bmatrix} + \begin{bmatrix} 0 \\ -7 \end{bmatrix}$$
        $$= 5\begin{bmatrix} 1 \\ 0 \end{bmatrix} - 7\begin{bmatrix} 0 \\ 1 \end{bmatrix}$$
        $$= 5\hat{i} - 7\hat{j}$$
    \end{exampleblock}
\end{frame}

% =============================================================================
\section{Linear Transformations and Matrices}
% =============================================================================

\begin{frame}{Matrix Multiplication as Composition}
    \begin{definition}[Linear Transformation and Composition]
        A \textbf{linear transformation} is a mapping that preserves vector addition and scalar multiplication (grid lines remain parallel/evenly spaced, origin is fixed). \textbf{Matrix composition} applies successive transformations via matrix multiplication.
    \end{definition}

    \begin{block}{Derivation}
        Applying a transformation matrix $A$ to a vector $\vec{x}$ is computed by taking the dot product of $A$'s rows with $\vec{x}$:
        $$\begin{bmatrix} a & b \\ c & d \end{bmatrix} \begin{bmatrix} x \\ y \end{bmatrix} = x\begin{bmatrix} a \\ c \end{bmatrix} + y\begin{bmatrix} b \\ d \end{bmatrix} = \begin{bmatrix} ax+by \\ cx+dy \end{bmatrix}$$
    \end{block}
\end{frame}

\begin{frame}{Matrix Multiplication as Composition (Example)}
    \begin{exampleblock}{Example Problem}
        \textbf{Problem:} Apply the shear transformation matrix $A = \begin{bmatrix} 1 & 2 \\ 0 & 1 \end{bmatrix}$ to the vector $\vec{v} = \begin{bmatrix} 1 \\ 2 \end{bmatrix}$. \\
        \vspace{0.5cm}
        \textbf{Solution:}
        $$A\vec{v} = \begin{bmatrix} 1 & 2 \\ 0 & 1 \end{bmatrix} \begin{bmatrix} 1 \\ 2 \end{bmatrix}$$
        $$= \begin{bmatrix} 1(1) + 2(2) \\ 0(1) + 1(2) \end{bmatrix}$$
        $$= \begin{bmatrix} 1 + 4 \\ 0 + 2 \end{bmatrix} = \begin{bmatrix} 5 \\ 2 \end{bmatrix}$$
    \end{exampleblock}
\end{frame}

\begin{frame}{Three-Dimensional and Non-square Matrices}
    \begin{definition}[Dimension Transformations]
        \textbf{Non-square matrices} represent linear transformations between spaces of differing dimensions. An $m \times n$ matrix transforms an $n$-dimensional space into an $m$-dimensional space.
    \end{definition}
    
    \begin{block}{Derivation}
        Mapping $T: \mathbb{R}^n \rightarrow \mathbb{R}^m$ uses an $m \times n$ matrix $A$. The product $A\vec{x}$ computes the coordinates in the target $m$-dimensional space, essentially taking a linear combination of the $n$ column vectors of $A$ (which live in $\mathbb{R}^m$).
    \end{block}
\end{frame}

\begin{frame}{Three-Dimensional and Non-square Matrices (Example)}
    \begin{exampleblock}{Example Problem}
        \textbf{Problem:} Map a 2D vector $\vec{x} = \begin{bmatrix} 2 \\ 3 \end{bmatrix}$ into 3D space using the transformation matrix $A = \begin{bmatrix} 1 & 0 \\ 0 & 1 \\ 1 & -1 \end{bmatrix}$. \\
        \vspace{0.5cm}
        \textbf{Solution:}
        $$A\vec{x} = \begin{bmatrix} 1 & 0 \\ 0 & 1 \\ 1 & -1 \end{bmatrix} \begin{bmatrix} 2 \\ 3 \end{bmatrix}$$
        $$= 2\begin{bmatrix} 1 \\ 0 \\ 1 \end{bmatrix} + 3\begin{bmatrix} 0 \\ 1 \\ -1 \end{bmatrix}$$
        $$= \begin{bmatrix} 2 \\ 0 \\ 2 \end{bmatrix} + \begin{bmatrix} 0 \\ 3 \\ -3 \end{bmatrix} = \begin{bmatrix} 2 \\ 3 \\ -1 \end{bmatrix}$$
    \end{exampleblock}
\end{frame}

% =============================================================================
\section{Determinants, Inverses, and Spaces}
% =============================================================================

\begin{frame}{The Determinant and Inverse Matrices}
    \begin{definition}[Determinant and Inverse]
        The \textbf{determinant} ($\det$) is the geometric scaling factor of a transformation's area or volume. The \textbf{inverse matrix} ($A^{-1}$) perfectly undoes the transformation $A$, such that $A^{-1}A = I$. An inverse only exists if $\det(A) \neq 0$.
    \end{definition}

    \begin{block}{Derivation for 2x2 Matrices}
        For a matrix $A = \begin{bmatrix} a & b \\ c & d \end{bmatrix}$: \\
        Determinant: $\det(A) = ad - bc$. \\
        Inverse: $A^{-1} = \frac{1}{\det(A)} \begin{bmatrix} d & -b \\ -c & a \end{bmatrix}$.
    \end{block}
\end{frame}

\begin{frame}{The Determinant and Inverse Matrices (Example)}
    \begin{exampleblock}{Example Problem}
        \textbf{Problem:} Find the inverse of $A = \begin{bmatrix} 4 & 3 \\ 3 & 2 \end{bmatrix}$. \\
        \vspace{0.5cm}
        \textbf{Solution:} \\
        1. Calculate the determinant:
        $$\det(A) = (4)(2) - (3)(3) = 8 - 9 = -1$$
        2. Apply the inverse formula:
        $$A^{-1} = \frac{1}{-1} \begin{bmatrix} 2 & -3 \\ -3 & 4 \end{bmatrix}$$
        $$A^{-1} = \begin{bmatrix} -2 & 3 \\ 3 & -4 \end{bmatrix}$$
    \end{exampleblock}
\end{frame}

\begin{frame}{Column Space, Null Space, and Cramer's Rule}
    \begin{definition}[Spaces and Cramer's Rule]
        The \textbf{column space} is the span of a matrix's columns. The \textbf{null space} contains all vectors mapped to the zero vector. \textbf{Cramer's Rule} is a theorem giving an explicit determinant-based solution for systems of linear equations $A\vec{x} = \vec{b}$.
    \end{definition}

    \begin{block}{Derivation of Cramer's Rule}
        To solve for the $i$-th variable $x_i$: substitute the $i$-th column of the coefficient matrix $A$ with the target vector $\vec{b}$ to form a new matrix $A_i$. 
        Then calculate: $$x_i = \frac{\det(A_i)}{\det(A)}$$
    \end{block}
\end{frame}

\begin{frame}{Cramer's Rule (Example)}
    \begin{exampleblock}{Example Problem}
        \textbf{Problem:} Solve the system $2x + y = 5$ and $x - y = 1$ using Cramer's rule. \\
        \vspace{0.2cm}
        \textbf{Solution:} 
        $A = \begin{bmatrix} 2 & 1 \\ 1 & -1 \end{bmatrix}$, $\vec{b} = \begin{bmatrix} 5 \\ 1 \end{bmatrix}$. $\det(A) = -2 - 1 = -3$. \\
        \vspace{0.2cm}
        Find $x$:
        $$A_x = \begin{bmatrix} 5 & 1 \\ 1 & -1 \end{bmatrix} \implies \det(A_x) = -5 - 1 = -6 \implies x = \frac{-6}{-3} = 2$$
        Find $y$:
        $$A_y = \begin{bmatrix} 2 & 5 \\ 1 & 1 \end{bmatrix} \implies \det(A_y) = 2 - 5 = -3 \implies y = \frac{-3}{-3} = 1$$
    \end{exampleblock}
\end{frame}

% =============================================================================
\section{Vector Products}
% =============================================================================

\begin{frame}{Dot Products and Duality}
    \begin{definition}[Dot Product and Duality]
        The \textbf{dot product} measures directional alignment between vectors. \textbf{Duality} highlights that any linear transformation mapping 2D vectors to a 1D number line exactly matches taking a dot product with a specific 2D vector.
    \end{definition}

    \begin{block}{Derivation}
        Algebraically, it is the sum of the products of corresponding entries: 
        $$\vec{u} \cdot \vec{v} = u_1 v_1 + u_2 v_2 + \dots + u_n v_n$$
        Geometrically, it relates to the angle $\theta$ between vectors: 
        $$\vec{u} \cdot \vec{v} = \|\vec{u}\| \|\vec{v}\| \cos(\theta)$$
    \end{block}
\end{frame}

\begin{frame}{Dot Products and Duality (Example)}
    \begin{exampleblock}{Example Problem}
        \textbf{Problem:} Calculate the dot product of $\vec{u} = \begin{bmatrix} 3 \\ -2 \\ 4 \end{bmatrix}$ and $\vec{v} = \begin{bmatrix} 1 \\ 5 \\ 2 \end{bmatrix}$. \\
        \vspace{0.5cm}
        \textbf{Solution:}
        $$\vec{u} \cdot \vec{v} = (3)(1) + (-2)(5) + (4)(2)$$
        $$= 3 - 10 + 8$$
        $$= 1$$
        \textit{Note: Since the dot product is positive, the angle between the vectors is acute.}
    \end{exampleblock}
\end{frame}

\begin{frame}{Cross Products and Linear Transformations}
    \begin{definition}[Cross Product]
        The \textbf{cross product} $\vec{u} \times \vec{v}$ produces a vector that is orthogonal (perpendicular) to both $\vec{u}$ and $\vec{v}$ in 3D space. Its magnitude is exactly equal to the area of the parallelogram spanned by $\vec{u}$ and $\vec{v}$.
    \end{definition}

    \begin{block}{Derivation via Determinant}
        The cross product is formally calculated using a pseudo-determinant of the basis vectors:
        $$\vec{u} \times \vec{v} = \begin{vmatrix} \hat{i} & \hat{j} & \hat{k} \\ u_1 & u_2 & u_3 \\ v_1 & v_2 & v_3 \end{vmatrix}$$
    \end{block}
\end{frame}

\begin{frame}{Cross Products and Linear Transformations (Example)}
    \begin{exampleblock}{Example Problem}
        \textbf{Problem:} Find $\vec{u} \times \vec{v}$ for $\vec{u} = \begin{bmatrix} 1 \\ 0 \\ 0 \end{bmatrix}$ (which is $\hat{i}$) and $\vec{v} = \begin{bmatrix} 0 \\ 1 \\ 0 \end{bmatrix}$ (which is $\hat{j}$). \\
        \vspace{0.5cm}
        \textbf{Solution:}
        $$\vec{u} \times \vec{v} = \begin{vmatrix} \hat{i} & \hat{j} & \hat{k} \\ 1 & 0 & 0 \\ 0 & 1 & 0 \end{vmatrix}$$
        $$= \hat{i}(0-0) - \hat{j}(0-0) + \hat{k}(1-0)$$
        $$= \hat{k} = \begin{bmatrix} 0 \\ 0 \\ 1 \end{bmatrix}$$
    \end{exampleblock}
\end{frame}

% =============================================================================
\section{Advanced Concepts}
% =============================================================================

\begin{frame}{Vector Spaces and Change of Basis}
    \begin{definition}[Vector Spaces and Change of Basis]
        A \textbf{vector space} is any set of objects obeying 8 strict axioms for vector addition and scalar multiplication. \textbf{Change of basis} translates vectors and transformations from our standard coordinate system into someone else's basis vectors.
    \end{definition}

    \begin{block}{Derivation}
        Let $P$ be a matrix whose columns are the new basis vectors. 
        \begin{itemize}
            \item To express a standard vector $\vec{v}$ in the new basis $P$, compute $[\vec{v}]_P = P^{-1}\vec{v}$.
            \item To translate a transformation matrix $A$ into the new basis, use $A' = P^{-1} A P$.
        \end{itemize}
    \end{block}
\end{frame}

\begin{frame}{Vector Spaces and Change of Basis (Example)}
    \begin{exampleblock}{Example Problem}
        \textbf{Problem:} Express the vector $\vec{v} = \begin{bmatrix} 3 \\ 2 \end{bmatrix}$ (in standard basis) in terms of a new basis defined by the columns of $P = \begin{bmatrix} 1 & 1 \\ 0 & 1 \end{bmatrix}$. \\
        \vspace{0.5cm}
        \textbf{Solution:} 
        1. Find $P^{-1}$:
        $$P^{-1} = \begin{bmatrix} 1 & -1 \\ 0 & 1 \end{bmatrix}$$
        2. Compute $[\vec{v}]_P = P^{-1}\vec{v}$:
        $$[\vec{v}]_P = \begin{bmatrix} 1 & -1 \\ 0 & 1 \end{bmatrix} \begin{bmatrix} 3 \\ 2 \end{bmatrix} = \begin{bmatrix} 3 - 2 \\ 0 + 2 \end{bmatrix} = \begin{bmatrix} 1 \\ 2 \end{bmatrix}$$
    \end{exampleblock}
\end{frame}

\begin{frame}{Eigenvectors and Eigenvalues}
    \begin{definition}[Eigenvectors and Eigenvalues]
        \textbf{Eigenvectors} ($\vec{v}$) are special, non-zero vectors that remain on their original span during a linear transformation $A$ (they only scale, they do not rotate). The scalar by which they stretch or compress is the \textbf{eigenvalue} ($\lambda$).
    \end{definition}

    \begin{block}{Derivation (Characteristic Equation)}
        The relationship is defined as $A \vec{v} = \lambda \vec{v}$, which rewrites to $(A - \lambda I)\vec{v} = \vec{0}$.
        To find non-trivial solutions for $\vec{v}$, we solve the characteristic equation:
        $$\det(A - \lambda I) = 0$$
    \end{block}
\end{frame}

\begin{frame}{Eigenvectors and Eigenvalues (Example - Part 1)}
    \begin{exampleblock}{Example Problem}
        \textbf{Problem:} Find the eigenvalues and corresponding eigenvectors of the matrix $A = \begin{bmatrix} 3 & 1 \\ 0 & 2 \end{bmatrix}$. \\
        \vspace{0.2cm}
        \textbf{Solution (Part 1 - Eigenvalues):} \\
        1. Set up the characteristic equation $\det(A - \lambda I) = 0$:
        $$\det\left( \begin{bmatrix} 3-\lambda & 1 \\ 0 & 2-\lambda \end{bmatrix} \right) = 0$$
        2. Calculate the determinant:
        $$(3-\lambda)(2-\lambda) - (1)(0) = 0$$
        $$(3-\lambda)(2-\lambda) = 0$$
        3. Solve for $\lambda$: The eigenvalues are $\lambda_1 = 3$ and $\lambda_2 = 2$.
    \end{exampleblock}
\end{frame}

\begin{frame}{Eigenvectors and Eigenvalues (Example - Part 2)}
    \begin{exampleblock}{Solution Continued (Part 2 - Eigenvectors)}
        \textbf{For $\lambda_1 = 3$:} Solve $(A - 3I)\vec{v} = \vec{0}$
        $$\begin{bmatrix} 3-3 & 1 \\ 0 & 2-3 \end{bmatrix} \begin{bmatrix} x \\ y \end{bmatrix} = \begin{bmatrix} 0 & 1 \\ 0 & -1 \end{bmatrix} \begin{bmatrix} x \\ y \end{bmatrix} = \begin{bmatrix} 0 \\ 0 \end{bmatrix}$$
        This yields $y = 0$, with $x$ as a free variable. Setting $x = 1$, the eigenvector is $\vec{v}_1 = \begin{bmatrix} 1 \\ 0 \end{bmatrix}$. \\
        \vspace{0.2cm}
        \textbf{For $\lambda_2 = 2$:} Solve $(A - 2I)\vec{v} = \vec{0}$
        $$\begin{bmatrix} 3-2 & 1 \\ 0 & 2-2 \end{bmatrix} \begin{bmatrix} x \\ y \end{bmatrix} = \begin{bmatrix} 1 & 1 \\ 0 & 0 \end{bmatrix} \begin{bmatrix} x \\ y \end{bmatrix} = \begin{bmatrix} 0 \\ 0 \end{bmatrix}$$
        This yields $x + y = 0 \implies x = -y$. Setting $y = 1$, the eigenvector is $\vec{v}_2 = \begin{bmatrix} -1 \\ 1 \end{bmatrix}$.
    \end{exampleblock}
\end{frame}

% =============================================================================
\section{References}
% =============================================================================
\begin{frame}[allowframebreaks]{References}
    \begin{thebibliography}{99}
        \setbeamertemplate{bibliography item}[text]
        
        \bibitem{bronson2023}
        Bronson, R., Costa, G.B., Saccoman, J.T. and Gross, D.
        \newblock \emph{Linear algebra: algorithms, applications, and techniques}. 4th ed., 2023.
        
        \bibitem{lehman2010}
        Eric Lehman, F Thomson Leighton, Albert R Meyer.
        \newblock \emph{Mathematics for Computer Science}. 1st ed., MIT, 2010.
        
        \bibitem{epp2010}
        Susanna S. Epp.
        \newblock \emph{Discrete Mathematics with Applications}. 4th ed., Brooks Cole, 2010.
        
        \bibitem{chartrand2012}
        Gary Chartrand, Ping Zhang.
        \newblock \emph{A First Course in Graph Theory}. 1st ed., Dover Publications, 2012.
        
        \bibitem{tsitsiklis2013}
        John Tsitsiklis.
        \newblock \emph{6.041SC Probabilistic Systems Analysis and Applied Probability}. Fall 2013. MIT OpenCourseWare. \url{https://ocw.mit.edu}.
        
        \bibitem{meyer2003}
        Albert Meyer.
        \newblock \emph{6.844 Computability Theory of and with Scheme}. Spring 2003. MIT OpenCourseWare. \url{https://ocw.mit.edu}.
        
        \bibitem{mitzenmacher2017}
        Michael Mitzenmacher and Eli Upfal.
        \newblock \emph{Probability and Computing}. 2nd ed., Cambridge University Press, 2017.
        
        \bibitem{granda2017}
        Carlos Fernandez-Granda.
        \newblock \emph{Course notes: DS-GA 1002: Probability and Statistics for Data Science}. \url{https://cims.nyu.edu/~cfgranda/pages/DSGA1002_fall17/index.html}.
        
    \end{thebibliography}
\end{frame}

\end{document}
